% ============================================================
%  Section 7 --- Discussion
% ============================================================

\section{Discussion}
\label{sec:discussion}

H1 est validée géométriquement (sil $= 0.519 \pm 0.092$, montée à $0.782 \pm 0.065$ en config
optimisée). H3 est validée \emph{au sens d'une discriminabilité statistique} (AUROC $= 0.973$
sur labels EWAT, $p < 0.01$ vs permutation A3, section~\ref{sec:robustness}). Cependant,
les stress tests A1 et A2 (section~\ref{sec:stress_tests}) montrent que cette discriminabilité
sur labels EWAT repose sur l'identification de scénarios connus depuis leur signature dans
l'espace latent, sur une cible \emph{circulaire} (les labels viennent du pipeline lui-même).

\textbf{L'évaluation indépendante sur cible Chaos Mesh} (section~\ref{sec:archv2}) donne
un \textbf{headline défensif de macro-AUROC $\mathbf{= 0.920}$ [0.878, 0.956] sur ewat\_v4\_strat}
avec LR-OvR sur features brutes flatten instance-normalisées (sans STGCN). Cette évaluation
est le bon chiffre à comparer à l'état de l'art : elle utilise une vérité terrain indépendante,
des IC bootstrap explicites, et un protocole stratifié-temporel.

La précursion temporelle est par ailleurs \textbf{confirmée} : $\Delta($far\,$-$\,near$) = -0.116$
sur le modèle STGCN+Chaos Mesh (Tableau~\ref{tab:c2_distant}) --- la dynamique pré-injection compte
pour 12\,pp d'AUROC.

H2a est réfutée définitivement, y compris sur les épisodes longs d'ewat\_v4\_strat (section~C5,
$p = 0.37$). H2b reste un PASS trivial.

Cette section discute ce que ces résultats apportent collectivement, comment ils se
positionnent par rapport à la baseline z-score et à l'état de l'art, puis liste
les limitations identifiées et les perspectives pour la prochaine itération.

\subsection{Ablation H1 vs.\ H3 --- géométrie $\neq$ prédictibilité}
\label{sec:ablation_h1_h3}

L'ablation par modalité révèle un résultat \emph{inverse} selon l'hypothèse testée.
Sur H1 (silhouette test, réentraînement complet), \texttt{M\_only} bat le modèle
\texttt{full} ($+0.058$ en sil\_test) : les features traces et logs ajoutent du bruit
à la géométrie siamoise sur $n=209$ épisodes train.
Sur H3 (AUROC précurseur, masquage à l'inférence), le modèle \texttt{full}
($0.954$) bat toutes les réductions ; \texttt{M\_only} tombe à $0.756$.
Les modalités T et L sont donc utiles pour la \emph{prédictibilité} des signatures
pré-injection, pas pour la structuration non supervisée des clusters.
Ce découplage doit être explicité dans toute interprétation des ablations feature-wise.

\subsection{Résultats négatifs comme contribution}
\label{sec:negative_results}

Un résultat négatif non exploitable est un échec : l'expérience n'apporte pas
d'information nouvelle. Un résultat négatif exploitable est une contribution :
il identifie précisément une contrainte, fournit une quantification et suggère
une direction corrective.

H2a appartient à la deuxième catégorie. Le mécanisme look-through MMD-RFF est
conçu correctement : il transmet le signal pendant le drift avec le flag approprié,
et le DriftDetector se recalibre lorsque le drift est stabilisé. Ce que la réfutation
de H2a révèle n'est pas un défaut de conception, mais une contrainte de données
quantifiée : sur des épisodes de $\approx 21$ steps, le warm-up du DriftDetector
($\approx 8$ steps) et la confirmation temporelle ($3$--$6$ steps) consomment jusqu'à
$14/21 = 67$\,\% de l'épisode avant que le mécanisme soit pleinement opérationnel.
Il ne reste alors que $\approx 7$ steps utiles --- une fenêtre trop courte pour que
le test de Student unilatéral atteigne $p < 0.05$ avec $n = 45$ épisodes.

La même contrainte explique H2b trivial : avec une fenêtre de confirmation de 5 steps,
le DriftDetector déclenche sur quasi tous les épisodes quel que soit leur type,
rendant l'overlap drift+alerte informatif pour aucun cluster en particulier.
L'OR de Fisher ($1.48$, $p = 0.35$) pour C8 vs drift pur confirme l'absence de
discrimination statistique.

Ces deux résultats négatifs sont reproductibles : le même protocole appliqué sur
ewat\_v4 avec des épisodes $\geq 40$ steps testera H2a dans des conditions où le
mécanisme peut effectivement opérer. Si H2a est validée sur ewat\_v4, la réfutation
sur ewat\_v3 aura permis d'identifier précisément la longueur minimale d'épisode
nécessaire.

\paragraph{L'alerte précède le drift flag.}
L'analyse temporelle de H2b révèle un résultat supplémentaire : l'alerte précurseur
précède le drift flag dans 85--100\,\% des épisodes (timing gap médian $< 0$
pour tous les clusters). Le DriftDetector est un indicateur tardif, pas précoce.
Cela est cohérent avec son fonctionnement : il mesure un changement de distribution
déjà établi, là où les précurseurs (étape 3) détectent des signatures
\emph{pré-injection}. Ce résultat renforce l'intérêt de l'étape 3 indépendamment
du mécanisme de détection de drift.

\subsection{Apport EWAT vis-à-vis du z-score}
\label{sec:ewat_vs_zscore}

La baseline z-score représente l'approche la plus répandue en production : une alarme
univariée sur chaque métrique, déclenchée lorsque la valeur dépasse $\sigma$ écarts-types
au-dessus de la moyenne glissante. Cette baseline est simple, explicable et robuste ---
mais elle souffre d'un défaut structurel : elle n'a aucun moyen de distinguer un drift
bénin d'une anomalie réelle, car les deux produisent des déviations statistiques similaires.
En conséquence, le taux de fausses alarmes sur les drifts reste à 100\,\% pour tout
$\sigma \in [2.0, 3.5]$.

EWAT au seuil 0.7 réduit ce taux à 8.3\,\% --- une réduction d'un facteur 12 --- tout
en conservant un lead time de 3.0\,min. Ce gain découle directement du typage contrastif :
les clusters de drift bénin (C5, C6, C9) ont des précurseurs entraînés sur leurs
signatures spécifiques, ce qui permet à l'AlertAssembler de les distinguer des anomalies.

Deux nuances sont importantes. D'abord, la détection d'anomalies réelles au seuil 0.7
est de 57.6\,\%, inférieure aux 100\,\% du z-score. Ce trade-off FA/détection est explicite
et contrôlable : le seuil 0.3 retrouve 100\,\% de détection avec 100\,\% de FA, le seuil 0.7
offre le meilleur équilibre sur les données disponibles. La courbe FA/détection n'était pas
disponible avec le z-score (le seuil $\sigma$ contrôle la sensibilité globale, pas la
discrimination drift/anomalie).
Ensuite, EWAT produit une sortie structurée --- le type $C_i$, la probabilité $\hat{p}_i$,
le lead time estimé $k^*_i$, et la fiche d'interprétabilité $\text{fiche}_{C_i}$ ---
là où le z-score ne produit qu'un binaire. Cette richesse informationnelle permet à
l'opérateur de prioriser et de préparer sa réponse avant la dégradation.

\subsection{Valeur de l'encodeur STGCN}
\label{sec:stgcn_value}

La comparaison B3/B4 (baselines précurseurs sur les scénarios Chaos Mesh) révèle que
l'encodeur STGCN n'améliore pas l'AUROC macro agrégé par rapport aux features brutes
(B3 $=$ B4 $= 0.835$). Cette égalité numérique est une coïncidence sur $n = 45$ épisodes
--- les AUROCs sont des rapports de paires concordantes entiers sur 126, et l'égalité
exacte résulte d'une redistribution de $\pm 75$ paires concordantes entre scénarios.

La valeur du STGCN est donc de deux natures. Géométriquement, il structure l'espace
latent de façon exploitable par le clustering siamois (H1 sil $= 0.519$ vs sil $\approx 0.3$
attendu sans structuration). Redistributivement, il améliore fortement certains types
difficiles (\texttt{fail\_slow\_cpu} $+0.270$, \texttt{drift\_scale\_up} $+0.111$)
au prix d'une dégradation sur des types déjà bien séparables dans l'espace brut
(\texttt{noisy\_neighbor} $-0.246$, \texttt{drift\_config\_change} $-0.127$).

Le résultat de l'ablation H1 (M\_only $+0.058$ vs full) indique que, sur 209 épisodes
d'entraînement, le signal des métriques Prometheus suffit à structurer l'espace latent
et que les traces et logs ajoutent du bruit géométrique. Ce résultat suggère que la valeur
de $T(t)$ et $L(t)$ est prédictive (H3) plutôt que géométrique (H1), ce qui justifie
de conserver les deux modalités dans le pipeline complet.

\subsection{Limitations}
\label{sec:limitations}

\begin{description}
  \item[\textbf{L1 --- Épisodes trop courts} ($\approx 21$ steps).]
    Conséquence directe sur H2a (look-through) et sur le plafond du lead time observable.
    L'AUROC précurseur est mesuré sur un horizon maximal $k^* = 10$ steps (5\,min),
    limité par la durée des épisodes. Le lead time réel pourrait être plus long si les
    précurseurs étaient visibles plus tôt dans des épisodes prolongés.
    Résolu dans ewat\_v4 ($\geq 40$ steps, $\geq 20$\,min par épisode).

  \item[\textbf{L2 --- disk\_io 16.7\,\% NaN} (\texttt{product-catalog}).]
    Le nœud \texttt{observit-cluster1-workers-58w74-mwxb2} est resté en état NotReady
    pendant l'intégralité de la collecte ewat\_v3. L'ablation LOO montre que
    \texttt{disk\_io} est la feature la plus critique pour H3 ($\Delta = -0.088$)
    malgré ses NaN --- son importance réelle est probablement sous-estimée.
    Résolu dans ewat\_v4 par déploiement d'un OTel SDK applicatif indépendant du nœud.

  \item[\textbf{L3 --- Méthode gradient invalidée}.]
    $\rho_\text{Spearman}(\text{gradient} \times \text{input},\, \text{permutation}) = -0.34$
    sur ewat\_v3 : la méthode gradient×input est anti-corrélée avec la permutation importance.
    Les fiches de clusters ont été régénérées avec \texttt{permutation\_importance}
    (50 shuffles), validées par KernelSHAP ($\rho > 0$ pour $9/10$ clusters
    \cite{lundberg2017}). Seul C3 (\texttt{noisy\_neighbor}) reste discordant ($\rho = -0.07$).
    Les fiches sont indicatives et non certifiées pour C3.

  \item[\textbf{L4 --- Transfert de domaine limité}.]
    Le pipeline entraîné sur ewat\_v3 détecte des anomalies génériques sur RCAEval
    (H1 sil $= 0.684$ avec instance norm + M\_only) mais ne discrimine pas les types
    (H3 AUROC $\approx 0.5$ pour tout $n_\text{few}$).
    La Stratégie A (refit scaler seul) est insuffisante.
    La Stratégie B (fine-tuning classifieurs LR ou dernière couche encodeur avec épisodes
    labellisés RCAEval) est planifiée pour ewat\_v4.

  \item[\textbf{L5 --- Transfer Entropy cluster-level : biais écologique}.]
    L'estimateur cluster-level agrège les séries temporelles de tous les épisodes d'un
    cluster avant d'estimer la TE, confondant les effets intra-épisode et inter-épisodes
    (biais écologique). L'estimateur service-level (intra-épisode) corrige ce biais
    et produit 124 relations causales significatives --- mais la TE reste univariée par
    feature, sous-estimant les synergies dans $\mathbb{R}^{17}$.

  \item[\textbf{L6 --- Faibles effectifs pour C6 et C9}.]
    $n_\text{pos\_test} = 1$ pour \texttt{drift\_config\_change} (C6) et
    \texttt{drift\_scale\_up} (C9) : l'AUROC H3 n'est pas calculable.
    Ces deux types de drift bénin requièrent un nombre d'épisodes test plus élevé,
    notamment $n_\text{pos\_test} \geq 5$ pour des IC bootstrap fiables.
\end{description}

\subsection{Perspectives --- ewat\_v4}
\label{sec:ewat_v4}

\paragraph{Données.}
La priorité principale est de corriger L1 et L2 simultanément. Les épisodes ewat\_v4
seront collectés sur $\geq 40$ steps ($\geq 20$\,min), laissant au DriftDetector un
budget de confirmation suffisant ($\leq 35$\,\% de l'épisode consommé par le warm-up
et la confirmation). Le nœud NotReady sera remplacé ou contourné par un OTel SDK
applicatif, ramenant le NaN de \texttt{disk\_io} à $0$\,\%.
Les scénarios C6 et C9 seront sur-représentés dans la collecte (objectif :
$n_\text{pos\_test} \geq 5$ par type) pour permettre le calcul des IC bootstrap H3.

\paragraph{Méthodologie.}
H2a sera réévaluée sur ewat\_v4 avec le même protocole (test de Student unilatéral
apparié, 45 épisodes test, seuil $p < 0.05$) pour tester si la longueur d'épisode
était bien le seul goulot d'étranglement. KernelSHAP sera réexécuté avec $n_\text{bg} \geq 50$
(vs 20 sur ewat\_v3) pour valider les fiches de clusters sur C3.
La Stratégie B du few-shot transfer sera implémentée : fine-tuning des classifieurs LR
sur des épisodes RCAEval labellisés, puis évaluation de H3 AUROC $> 0.7$ en transfert.

\paragraph{Architecture.}
Le GAT est le candidat naturel pour ewat\_v4 : il produit la meilleure géométrie
(sil\_test $= 0.497$) et le plus grand nombre de types prédictibles ($13/15$).
Sur des épisodes plus longs, le mécanisme d'attention adaptative sur $G(t)$ devrait
amplifier son avantage par rapport au STGCN.
Les features à $\Delta \approx 0$ en LOO (\texttt{retry\_rate}, \texttt{log\_warn\_rate},
\texttt{queue\_depth}) seront supprimées, réduisant $S(t)$ de $\mathbb{R}^{N \times 17}$
à $\mathbb{R}^{N \times 14}$ --- ce qui simplifiera l'encodeur sans perte mesurable
de prédictibilité.
La TE multivariée dans $\mathbb{R}^{17}$ (estimateur JIDT ou k-NN multivarié)
remplacera la somme des TE univariés pour capturer les synergies entre features,
corrigeant la limitation L5.

\subsection{Perspectives --- au-delà d'ewat\_v4 (post-audit C-1 à C-5)}
\label{sec:future_work}

L'audit post-Phase~C (section~\ref{sec:residual_limits} et
\texttt{docs/limitations.md} L10--L17) a identifié huit limites résiduelles qui
orientent la suite du programme de recherche EWAT :

\begin{enumerate}
  \item \textbf{Stabilisation H1 sur datasets larges} (L10) : implémenter
    le hard-negative mining renforcé pour réduire le surentraînement du siamois.
    Cible : sil\_test $\geq 0.60$ sur 10 graines de ewat\_v4.
  \item \textbf{Open-set state-of-the-art} (L12) : étendre
    \texttt{src/ewat/openset/} avec Mahalanobis-OOD (Lee 2018), Energy-based
    (Liu 2020), ODIN (Liang 2017). Cible : Unknown AUROC $\geq 0.7$ sur LOSO
    v4\_strat (vs $0.55$ actuel avec OpenMax).
  \item \textbf{Validation cross-cluster} (L14) : porter le pipeline sur un
    second cluster Kubernetes (EKS, GKE, ou cluster Devoteam différent) avec
    Stratégie B fine-tuning sur quelques épisodes du cluster cible.
  \item \textbf{Scaling à $N \geq 20$ services} (L13) : collecter des épisodes
    sur un benchmark microservices plus large (e.g., Train-Ticket, Hipster Shop
    étendu, ou un déploiement production réel). Re-benchmarquer la latence.
  \item \textbf{Audit ontologique SRE} (L17) : faire valider l'ontologie OWL/RDF
    Phase 8 par un SRE staff. Intégrer l'export RDF avec un système de
    runbook/alerting (PagerDuty rules, OpsGenie policies).
  \item \textbf{Cycle de retraining opérationnel} (L15) : implémenter un
    moniteur de drift d'embeddings (production vs entraînement) qui déclenche
    automatiquement un retraining quand un seuil est dépassé.
  \item \textbf{Auto-feature-engineering} (L16) : remplacer les 17 features
    fixées par une sélection automatique (mRMR ou SHAP-based pruning) qui
    s'adapte à chaque déploiement.
  \item \textbf{Validation production} : déployer EWAT en mode shadow
    (observation sans action) sur un cluster production réel pendant
    plusieurs semaines pour observer les vrais incidents et comparer aux
    chaos-injectés.
\end{enumerate}

L'élément critique pour la défense de cette itération reste néanmoins le
\textbf{headline défensif honnête} (macro-AUROC $= 0.920$ [0.878, 0.956] sur
cible Chaos Mesh indépendante, ewat\_v4\_strat, LR-OvR sur features brutes
instance-normalisées) avec son benchmark de latence vérifié (TOTAL p95
$= 13.28$ ms, 375$\times$ sous budget) et son audit power statistique
(5/10 clusters reportables avec power $= 1.0$ sur les reportables). Ces trois
résultats convergent vers une contribution prédictive crédible et opérationnellement
viable, accompagnée d'une documentation transparente de ses limites.
